% chapters/chapter1/sec2.tex

\section{The irreplaceable core value of ISAC}
\label{sec:ch1-isac-core value}

% TODO: 本节由 Huixin Man 负责撰写。
% 写作目标：
% 1. 给出 ISAC 的基本定义；
% 2. 说明通信与感知融合的主要形式；
% 3. 简要介绍 ISAC 的发展现状和典型应用场景。

上一节从概念层面回答了"ISAC是什么"以及"ISAC如何演进"的问题。本节将进一步回答一个更为根本的问题：ISAC的核心价值究竟是什么？为什么说它是不可替代的？
我们将从频谱效率、硬件成本、部署效率和业务闭环能力四个维度展开分析，揭示ISAC相较于传统独立系统所具备的底层优势。更重要的是，我们将说明这些优势并非孤立存在，而是共同构成了一种系统级的价值重构——正是这种重构，使得ISAC从一个"锦上添花"的技术选项，上升为6G网络不可或缺的内生能力。

\subsection{Why isn't ISAC merely an "optional addition to existing functions"?}
\label{subsec:ch1-isac-core value-why}
\subsubsection{The value logic from independent systems to integrated systems}
\label{subsubsec:ch1-isac-core value-logic}

在过去数十年中，通信系统和感知系统各自形成了完整且成熟的技术体系，各自解决着界限分明的问题。通信系统的全部努力都指向一个核心目标：将信息从一个节点可靠、高效地传送到另一个节点。无论是从1G到5G的代际演进，还是从FDMA到OFDMA的多址技术变迁，通信系统的设计始终围绕着信道容量、频谱效率、传输可靠性等指标展开。通信系统对外部物理世界的态度是"不关心"的——它关心的是比特的传输，而非比特所描述的物理现实。与之对应，独立的感知系统（以雷达为代表）则专注于对环境的探测和理解。它关心的是目标在哪里、以什么速度运动、环境中存在什么障碍物，但它并不负责将这些信息传递给远端的决策者或执行者。感知系统获取的信息如何被利用、被谁利用、如何与其他系统协同，这些问题长期以来不在感知系统自身的设计范畴之内。

这种分工在早期阶段是合理且高效的。然而，当我们审视面向6G的未来网络业务需求时，会发现一个根本性的变化：越来越多的关键应用不再满足于单纯的"信息传输"或单纯的"环境探测"，而是要求系统同时具备信息传输能力、环境理解能力和实时反馈能力，并将这三种能力紧密联动，形成完整的服务链条。以自动驾驶为例，车辆需要实时感知周围的车辆、行人和障碍物（环境理解），将感知结果与来自路侧基础设施和其他车辆的信息进行融合（信息传输），并据此做出驾驶决策和执行控制动作（实时反馈）。在这个过程中，感知和通信不是两个可以分别完成再简单拼合的独立任务，而是一个有机整体中相互依存的环节。感知的结果影响着通信的内容和优先级，通信传递的信息又丰富着感知的视野和精度，二者共同支撑着最终的决策质量。

ISAC的核心价值正在于此：它打通了感知、通信与业务决策之间的链条，使得无线网络从一个单纯的"信息管道"升级为一个能够"感知世界、传递信息、辅助决策"的综合服务平台。这种价值不是通过在通信系统旁边加装一套独立的感知设备就能实现的，因为松耦合的系统拼接无法提供紧密联动所需的实时性、一致性和效率。

\subsubsection{The core value of system-level revenue restructuring}
\label{subsubsec:ch1-isac-core value-revenue}

理解ISAC的价值时，一个常见的误区是将其简化为某个单一性能指标的提升——例如，认为ISAC的价值在于"感知精度超过独立雷达"或"通信速率超过独立通信系统"。这种理解是片面的，甚至在很多情况下是不成立的。事实上，如果仅从单一功能的角度比较，一套专门设计的独立雷达系统在感知性能上很可能优于ISAC系统中的感知功能，因为独立雷达可以将全部资源（功率、带宽、处理能力）集中用于感知任务，不需要与通信功能分享。同样，一套专用的通信系统在传输性能上也可能优于ISAC系统中的通信功能。这就好比一位专业的短跑运动员在百米赛跑中几乎必然快于一位十项全能选手，但后者的价值不在于任何单项的极致成绩，而在于综合能力的均衡与协同。

ISAC的真正价值体现为系统级的收益重构，具体表现为以下几个方面的综合提升：

\begin{enumerate}[label=(\arabic*)]
\item 资源利用率的整体提升。同一块频谱、同一套硬件、同一个时隙被同时用于通信和感知，系统整体的资源利用率显著高于两套独立系统的资源利用率之和。

\item 系统总成本的显著降低。通过硬件共享和基础设施复用，ISAC系统的建设成本和运维成本远低于分别部署通信系统和感知系统的总成本。

\item 部署灵活性和扩展性的增强。ISAC可以依托现有的通信网络基础设施快速引入感知能力，无需额外建设独立的感知基础设施，大大缩短了新服务的上线周期。

\item 服务能力的质变。ISAC使得网络具备了"感知—通信—决策—反馈"的闭环能力，这是独立的通信系统或独立的感知系统都无法单独提供的。这种闭环能力为上层应用提供了全新的服务范式。
\end{enumerate}

因此，评价ISAC的价值不能仅用通信领域或雷达领域的传统单一指标来衡量，而需要从系统整体的维度，综合考虑频谱效率、硬件成本、部署效率和业务闭环能力等多个方面。以下四个小节将逐一深入分析这些维度的优势.

\subsection{Spectrum efficiency advantage}
\label{subsec:ch1-isac-core value-spectrum}
\subsubsection{The spectrum occupation issue of traditional independent systems}
\label{subsubsec:ch1-isac-core value-spectrum-occupation}

频谱资源是无线系统运行的物理基础，也是当今最为稀缺的战略性资源之一。在传统的频谱管理体制下，通信系统和雷达系统各自拥有独立划分的工作频段，二者之间通过频率隔离来避免互干扰。这种管理模式虽然简单清晰，但带来了日益严重的频谱利用效率问题。

从全球频谱分配的总体格局来看，军用和民用雷达系统占用了相当可观的频谱资源，特别是在L波段（1–2 GHz）、S波段（2–4 GHz）、C波段（4–8 GHz）和X波段（8–12 GHz）等频率范围内。与此同时，移动通信系统对频谱的需求随着数据流量的指数增长而持续膨胀。根据行业预测，到2030年全球移动数据流量将达到2020年的100倍以上，这意味着通信系统需要获取远超当前分配量的频谱资源。
然而，这两类系统对各自频段的实际使用率往往并不高。雷达系统，特别是军用雷达和气象雷达，通常具有间歇式工作特征——发射脉冲之间存在大量的空闲时间，在频率维度上也并非持续占满整个分配频段。由于频谱管理的刚性划分，这些空闲的时频资源无法被通信系统利用。反之亦然：通信系统在某些时间和区域的频谱利用率也可能很低（如深夜的城区基站或偏远地区的基站），但这些资源同样无法被感知任务利用。
这种"各自占用、互不共享"的模式导致了严重的频谱碎片化问题。从全局视角来看，大量的频谱资源在时间和空间维度上处于低效利用状态，而两类系统却都在抱怨频谱不够用。这就好比一座城市中，住宅用地和商业用地被严格隔离，住宅区白天大量空置、商业区夜间无人使用，但城市却面临着"用地紧张"的困境。

随着通信系统向毫米波和太赫兹频段扩展，频谱冲突问题变得更加尖锐。在24–30 GHz和60–80 GHz等频段，通信业务与雷达业务（如汽车雷达、气象雷达、航空雷达）的频谱需求高度重叠，传统的频率隔离策略已经难以为继。

\subsubsection{The sources of ISAC's advantages in spectrum sharing and reuse}
\label{subsubsec:ch1-adv source of isac-spectrum}

ISAC对频谱效率的提升，不是简单地让通信和感知"挤在同一个频段里"，而是通过多个层面的机制实现频谱资源的深度复用。
最直接的频谱效率增益来源于同一时频资源的功能复用。在ISAC系统中，同一个OFDM符号或同一个发射波束可以同时承载通信数据和感知功能。通信接收端从中解调出数据信息，感知接收端则通过分析该信号的回波提取目标参数。从频谱使用的角度看，原来需要两块独立频谱分别完成的通信任务和感知任务，现在在同一块频谱上同时完成了。

如果我们用一个简化的例子来理解：假设通信系统需要100 MHz带宽来传输数据，雷达系统需要100 MHz带宽来实现所需的距离分辨率。在传统的独立部署模式下，两个系统共计需要200 MHz频谱。而在ISAC系统中，同一个100 MHz带宽的信号同时服务于通信和感知，总频谱需求降低了一半。当然，实际系统中的情况远比这个例子复杂，因为通信和感知对频谱的需求特征并不完全相同，联合使用时可能存在性能折中,但基本的频谱复用增益是确实存在的。
更深层次的频谱效率增益来源于联合信号设计。在ISAC系统中，信号波形不是分别针对通信和感知独立设计后再简单叠加，而是在设计之初就同时考虑两种功能的需求。联合设计可以巧妙地利用信号的不同维度（如时域、频域、空域、码域）来同时满足通信和感知的需求，从而在给定的频谱资源上获取超越简单"一分为二"的性能。例如，通过精心设计OFDM信号的子载波分配和导频图案，可以在保障通信解调性能的同时，为感知提供足够的参考信号资源，而这些导频信号本身就是通信系统正常运行所必需的,即感知功能在一定程度上"免费"搭乘了通信系统的信号结构。

在网络层面，ISAC还可以实现通信与感知任务的频谱协同调度。蜂窝网络中的不同基站、不同时隙面临的通信负载和感知需求是动态变化的。通过网络级的智能调度，可以在通信负载较轻的时隙和频段上增强感知功能，在感知需求较低的时段将更多资源分配给通信。这种动态的频谱协同调度在独立系统中是无法实现的，因为独立系统之间缺乏高效的信息交互和协调机制。

\subsubsection{The applicable conditions and boundaries of the spectral efficiency advantage}
\label{subsubsec:ch1-isac-core value-spectrum-boundary}

在充分认识ISAC频谱效率优势的同时，我们必须避免过度理想化的认识。频谱效率增益并非在所有场景和条件下都能自动获得，它受到多种因素的约束和限制。

首先，通信与感知对信号的需求存在内在矛盾，这种矛盾在某些场景下可能显著限制频谱复用的增益。通信功能希望信号承载尽可能多的随机数据符号以最大化信息传输速率，而感知功能（特别是距离-多普勒估计）则受益于信号的确定性和良好的相关特性。当同一个信号需要同时服务于两种功能时，必然存在性能折中——为通信性能分配更多的资源（如更多的数据子载波）可能导致感知性能下降，反之亦然。在某些极端场景下，通信和感知对频谱资源的需求差异非常大（例如，通信需要极高的频谱效率而感知需要极高的距离分辨率），联合使用同一频谱的增益可能被严重的性能折中所侵蚀。
其次，自干扰问题是ISAC系统频谱效率面临的一个独特挑战。在单站ISAC系统中，基站在发射通信/感知信号的同时需要接收感知回波。由于收发同时进行且使用相同的频率，发射信号会对接收机产生强烈的自干扰。如果自干扰不能被有效抑制，感知信号的信噪比将严重恶化，实际可获得的频谱效率增益将大打折扣。
再次，感知目标的参数需求也约束着频谱效率增益的大小。高精度的距离估计需要大信号带宽，高精度的速度估计需要长观测时间，高角度分辨率需要大天线孔径。这些感知需求可能与通信系统的资源分配策略产生冲突。例如，如果通信系统为了服务更多用户而采用短时隙设计，那么感知的速度估计精度可能受到限制。
最后，多用户多目标场景中的复杂干扰关系也可能制约频谱效率。当系统需要同时服务多个通信用户和检测多个感知目标时，通信用户之间、感知目标之间以及通信与感知之间的干扰管理变得更加复杂，简单的频谱复用可能无法获得理想的增益。

因此，对ISAC频谱效率优势的理解应该是条件化的：在合理的系统设计和适宜的应用场景下，ISAC能够带来显著的频谱效率提升；但这种提升的程度取决于具体的场景需求、系统参数和设计方案。后续章节在讨论具体的波形设计和资源优化时，将对这些条件和边界进行定量分析.

\subsection{Hardware cost advantage}
\label{subsec:ch1-isac-core value-hardware}
\subsubsection{Independent deployment of hardware duplication issue}
\label{subsubsec:ch1-isac-core value-hardware-duplication}

如果通信和感知功能分别由独立的系统承担，那么两套系统各自需要完整的硬件链路——从天线阵列、射频前端（包括功率放大器、低噪声放大器、混频器、滤波器等）、模数/数模转换器、基带信号处理单元，到时钟源、电源模块，乃至安装平台（如铁塔、机柜、散热系统等）。这种硬件的大面积重复建设带来了显著的经济负担。

以城市道路的车路协同场景为例。如果采用独立部署模式，每个路侧节点需要同时安装一套通信设备（如C-V2X路侧单元）和一套感知设备（如毫米波雷达或激光雷达）。两套设备各自拥有天线、射频链路和处理器，各自需要独立的供电和安装支架，各自需要独立的回传链路将数据传回边缘或云端。按照一个中等城市数千个路口计算，硬件的重复投资是非常可观的。
除了一次性的建设成本（CAPEX）之外，运维成本（OPEX）同样不容忽视。两套独立设备意味着双倍的维护工作量、双倍的能耗、双倍的故障排查复杂度。特别是在大规模部署的场景下，运维成本往往占据系统全生命周期成本的主要部分.

\subsubsection{The basic logic of ISAC hardware sharing}
\label{subsubsec:ch1-isac-core value-hardware-sharing}

ISAC的硬件成本优势来源于在多个层面上实现通信与感知功能的硬件共享。
在前端层面，通信与感知可以共享天线阵列和射频链路。在毫米波和太赫兹频段，通信系统和感知系统都需要使用相控阵天线来获得足够的波束增益以克服路径损耗。如果采用ISAC设计，同一套相控阵天线既可以生成指向通信用户的数据波束，也可以生成用于探测感知目标的扫描波束。射频链路中的功率放大器、低噪声放大器、混频器等组件同样可以被两种功能共享。基带处理单元中的FFT/IFFT运算模块、信道估计模块等，在通信和感知处理中也存在大量的可复用计算结构。

在平台层面，通信基站、路侧单元（RSU）、边缘计算节点、车载通信设备等现有平台可以直接作为ISAC的承载平台，无需另行建设独立的感知基础设施。以蜂窝基站为例，一座5G基站已经具备了大规模MIMO天线阵列、宽带射频前端、高性能基带处理器和稳定的回传链路——这些恰恰也是实现高质量雷达式感知所需要的核心硬件能力。通过软件升级和适度的硬件增强，现有基站就可以在执行通信任务的同时提供感知服务，边际硬件成本远低于独立部署一套专用感知设备。

在软件功能层面，通信和感知之间也存在大量可复用的信号处理模块。例如，通信系统中的同步模块（用于时间和频率同步）可以被感知功能复用来实现对目标回波的精确时间对齐；通信系统中的信道估计模块与感知中的目标参数估计在数学结构上具有相似性，可以共享部分算法框架；通信系统中的波束训练和波束管理过程本身就蕴含了对空间环境的扫描和探测，其结果可以直接用于感知.

\subsubsection{The limitations of hardware sharing and the engineering costs}
\label{subsubsec:ch1-isac-core value-hardware-limitations}

硬件共享带来了显著的成本优势，但共享并非没有代价。通信功能和感知功能对硬件的部分性能指标存在不同甚至矛盾的需求，这些差异给硬件设计带来了额外的挑战和成本。

\begin{itemize}
\item\textbf{动态范围:} 在雷达式感知中，发射信号经过目标反射后到达接收机时，功率可能比发射功率低100 dB以上（取决于目标的距离和雷达截面积）。这要求接收机具有非常高的动态范围，以便同时处理强近场回波和弱远场回波。而通信接收机虽然也需要一定的动态范围来适应不同距离用户的信号强度差异，但其动态范围需求通常低于雷达接收机。如果硬件平台需要同时满足感知对高动态范围的需求和通信对成本的敏感性，设计上需要找到合适的平衡点。

\item\textbf{收发隔离:} 在全双工ISAC系统中（即基站同时发射和接收），发射信号对接收机的泄漏（自干扰）可能高达发射功率本身的水平，而接收感知回波的信号功率可能比自干扰低80–100 dB。实现如此高的收发隔离度需要在天线设计、射频电路和数字信号处理等多个层面采取复杂的消除措施，这会增加硬件复杂度和成本。

\item\textbf{时钟精度和相位噪声:} 高精度的速度估计要求系统具有极低的相位噪声水平，因为目标的微多普勒频移可能非常小。通信系统对相位噪声的容忍度通常较高（因为通信系统通过信道估计和均衡来补偿相位偏转），但感知系统对相位噪声更为敏感。共享时钟和本振时，需要按照感知的更高精度要求来设计，这可能增加器件成本。

\item\textbf{功率放大器的线性度:} 通信系统中广泛使用的OFDM信号具有较高的峰均功率比（PAPR）,对功率放大器的线性度要求较高。而传统雷达信号（如恒模的LFM脉冲）对线性度的要求相对宽松，但对功率效率的要求更高。ISAC系统需要兼顾线性度和效率，这对功率放大器的设计提出了更高的要求。
\end{itemize}

因此，ISAC的硬件共享不是"免费的午餐"，而是通过增加一定的设计复杂度来换取整体成本的降低。从工程经济学的角度看，只要共享带来的成本节省大于为满足双功能需求而增加的设计成本，硬件共享就是有利的。具体的硬件非理想性分析和工程实现挑战将在后续相关章节中详细讨论.

\subsection{Deployment efficiency advantage}
\label{subsec:ch1-isac-core value-deployment}
\subsubsection{Why does the future network require high deployment efficiency}
\label{subsubsec:ch1-isac-core value-deployment-why}

未来网络的发展正在从传统的“面向通信连接”逐步走向“通信、感知、计算、控制”深度融合的新阶段。随着数字经济和智能社会的持续推进，网络服务对象、覆盖场景和业务需求都在快速扩张，这使得网络基础设施不仅要具备更强的性能，还必须具备更高的部署效率。

首先，新兴应用场景正在快速涌现并大规模扩张。例如，在低空经济领域，无人机物流、巡检、安防和应急救援等业务需要网络同时提供稳定通信与空域感知能力；在车路协同场景中，智能网联汽车需要道路基础设施支持环境感知、目标检测、轨迹预测与低时延信息交互；在工业互联网场景中，生产设备、移动机器人和自动化系统要求网络实现高可靠连接与生产环境实时感知；在公共安全场景中，重大活动保障、灾害监测、重点区域安防等任务也对网络的快速建网和态势感知能力提出了更高要求。上述场景往往覆盖范围广、变化速度快、部署地点复杂，如果网络建设周期过长，将难以及时支撑业务落地和规模化推广。

其次，若感知能力仍然主要依赖独立建设的专用基础设施，将显著增加系统部署复杂度与建设周期。传统模式下，通信系统和感知系统通常分别设计、分别部署：通信依赖基站、接入点等网络设施，而感知则依赖雷达、摄像头、激光雷达或其他专用传感设备。这种“烟囱式”建设方式虽然在单一功能优化上具有一定优势，但在未来大规模场景中会带来明显问题。一方面，独立部署感知基础设施意味着需要额外的站址、供电、回传链路、维护系统和频谱协调机制，导致建设成本和运维复杂度大幅上升；另一方面，多系统并行建设还可能造成资源重复投入、接口不统一、协同困难等问题，从而延长部署周期，降低网络建设效率。

再次，未来网络面向的应用环境具有显著的动态性和不确定性，要求网络具备快速开通、弹性扩展和按需升级的能力。例如，临时活动保障、灾后应急通信、阶段性工业产线调整、临时低空航线开放等，都需要网络在较短时间内完成覆盖部署和能力加载。如果每增加一种能力都需要建设一套新的专用设备体系，那么网络将难以适应这种快速变化的业务需求。因此，高部署效率不仅是降低成本的要求，更是提升网络适应性和服务敏捷性的关键条件。

从网络演进角度看，高部署效率已经成为未来网络基础能力的重要组成部分。它不仅体现在设备安装速度快、开通周期短，还体现在网络能力能够复用已有基础设施、通过软件定义和功能融合实现快速上线。特别是在通信与感知一体化（ISAC）框架下，如果能够基于现有通信网络设施扩展感知能力，就有望避免重复建设，缩短部署周期，实现“建网即具备感知能力”或“在网快速升级感知能力”的目标。这种模式能够更好地满足未来多场景、广覆盖、低成本、快响应的网络发展需求。

因此，未来网络之所以需要高部署效率，根本原因在于：应用场景扩张速度快、业务需求变化频繁、传统独立建网方式成本高且周期长，而网络能力融合化、软件化和平台化的发展趋势又为高效部署提供了技术基础。在这一背景下，ISAC所强调的“通信基础设施兼顾感知功能”的思路，正成为提升未来网络部署效率的重要方向。

\subsubsection{ISAC's mechanism for enhancing deployment efficiency}
\label{subsubsec:ch1-isac-core value-deployment-mechanism}

ISAC大幅提升部署效率的核心机制在于：依托已有的通信网络基础设施快速扩展感知能力。
当前在全球范围内，蜂窝通信网络已经建立了一个覆盖广泛、站点密集的基础设施网络。以中国为例，截至2024年底，全国5G基站数量已超过400万座，覆盖了绝大多数的城市、乡镇和主要交通干线。这些基站各自拥有高性能的天线阵列、射频前端和信号处理能力，还配备了完善的供电系统、回传网络和网管系统。如果能够通过软件升级和适度的硬件增强，使这些基站在执行通信任务的同时提供感知服务，那么就相当于"一夜之间"在全国范围内建立了一个拥有数百万节点的感知网络——而这一过程不需要新建站址、不需要新增回传链路、不需要新建网管平台，部署效率的提升是数量级的。

除了利用现有基础设施之外，ISAC还在新建站点的规划和建设中带来了效率提升。当通信和感知功能由同一个平台承载时，站点规划只需要进行一次（而非通信和感知分别规划），供电系统只需要建设一套，回传链路只需要一条，边缘处理平台只需要一个。这种统一规划的方式不仅降低了单站的建设成本，还简化了网络管理的复杂度。

从新业务上线的角度看，ISAC也显著降低了门槛。在传统模式下，如果一个城市希望开展低空无人机管控服务，需要从零开始规划建设一套独立的探测感知网络，整个过程涉及选址、环评、招标、施工、调试等多个环节。而在ISAC模式下，运营商只需要对现有基站进行软件升级或有限的硬件扩展，就可以快速上线感知服务，将业务部署周期从"年"级别缩短到"月"甚至"周"级别。

\subsubsection{From single-node efficiency to network-level efficiency}
\label{subsubsec:ch1-isac-core value-deployment-network}

ISAC的部署效率优势不仅体现在单个节点的层面，更体现在网络级的规模效应上。
在单站层面，一个ISAC基站替代了原来需要分别部署的通信基站和感知设备，安装、调试和维护的工作量显著减少。更重要的是，由于通信和感知共享同一套硬件和回传链路，站点的体积、重量和功耗也相应降低，这对于铁塔承载能力有限、机房空间紧张的城市密集部署场景尤为重要。

在网络层面，ISAC的优势更为显著。蜂窝通信网络已经建立了完善的多站协同机制，包括站间同步、干扰协调、切换管理和负载均衡等。当感知功能内生于蜂窝网络之中时，多站协同感知网络可以自然地继承这些已有的协同机制，而不需要从零开始构建感知网络的协同框架。例如，多基站联合感知需要各基站之间的精确时间同步——而蜂窝网络中的基站通常已经通过GPS或IEEE 1588精确时间协议实现了纳秒级的时间同步，这一基础设施直接可以被多站协同感知所利用。

从网络演进的角度看，ISAC使得网络升级路径更加平滑。运营商可以先在部分热点区域的基站上开通感知功能，根据业务需求和运营经验逐步扩展到更多站点，最终实现全网范围的通感一体化覆盖。这种渐进式的升级路径远比一次性大规模部署独立感知网络更加可控、风险更低.

\subsection{Advantage in Business Closed-Loop Capability}
\label{subsec:ch1-isac-core value-closed loop}

与传统网络主要承担“信息传递”功能不同，通信感知一体化（ISAC）不仅能够完成数据传输，还能够通过对外部环境、目标状态和业务对象的实时感知，形成从信息获取到业务执行再到结果反馈的完整链路。这样的能力使网络从“连接工具”进一步演进为“业务闭环支撑平台”，从而更好地适配未来智能化、自动化和实时化应用需求。
业务闭环能力的核心在于：网络不再只是被动承载业务数据，而是能够参与环境认知、状态判断、任务决策和执行反馈等关键环节。对于6G时代面向智能交通、低空经济、工业自动化和公共安全等场景的网络而言，这种能力将成为网络体系升级的重要方向。

\subsubsection{The limitations of traditional communication systems}
\label{subsubsec:ch1-isac-core value-closed loop-communication}
传统通信系统的主要目标是实现信息在用户、终端和业务平台之间的可靠传递，其设计重点通常集中在带宽、时延、可靠性、连接数和覆盖能力等指标上。从这一意义上看，传统网络在“传递信息”方面已经形成了较为成熟的体系，但其对外部物理环境本身往往缺乏直接认知能力。

具体而言，传统通信系统虽然能够承载来自传感器、终端设备和应用平台的数据流，但网络本身通常并不能直接获取外部环境状态，也难以主动识别目标位置、运动轨迹、行为变化或空间分布等信息。换言之，传统网络更擅长“把信息送到哪里”，却不具备天然的“知道环境中发生了什么”的能力。这种局限决定了其在面向智能业务时，通常只能作为信息通道存在，而难以直接参与业务状态理解与实时控制。
进一步来看，传统通信系统中的决策过程往往依赖于外部传感器系统或离线数据源。例如，在车路协同中，环境感知往往依赖摄像头、毫米波雷达、激光雷达等独立设备；在工业自动化中，状态检测常常依赖专用传感器网络；在公共安全场景中，异常识别和事件判断通常依赖视频监控系统或历史数据库。通信网络本身仅负责将这些感知数据上传或下发控制信息，而真正的分析与决策能力分散在网络外部的独立系统中。

这种模式带来了几个明显问题：一是系统链条较长，感知、传输、决策和执行之间存在较多环节，容易引入额外时延；二是多系统之间接口复杂、协同困难，影响整体响应效率；三是网络难以形成对业务过程的全局认知和统一调度能力。因此，传统通信系统在支撑未来高实时性、高自治性业务方面存在天然局限。

\subsubsection{ISAC forms a closed loop of "perception $-$ transmission $-$ decision $-$ feedback"}
\label{subsubsec:ch1-isac-core value-closed loop-isac}

ISAC的重要价值之一，在于它能够推动网络形成“感知—传输—决策—反馈”的业务闭环，使网络从单纯的信息承载平台演进为能够参与业务运行过程的智能基础设施。

\begin{enumerate}[label=(\arabic*)]
\item\textbf{感知获取环境状态}

通过一体化的信号设计、波形复用和网络基础设施协同，ISAC能够在通信过程中同步感知目标的位置、距离、速度、方向以及周边环境变化等信息。与传统模式下依赖完全独立的专用感知设备不同，ISAC强调利用通信系统已有资源获取对外部物理世界的实时认知能力，从而使网络能够及时掌握业务运行所依赖的环境状态。

\item\textbf{通信完成状态共享与控制信息下发}

感知获得的信息需要在网络内快速传输并共享给相关节点、平台或应用系统，包括边缘节点、控制中心、车载终端、工业控制器和空域管理平台等。同时，基于决策结果生成的控制指令，也需要通过通信链路及时、可靠地下发到执行对象。由此，通信在业务闭环中不仅承担数据传输功能，还承担状态同步、协同控制和任务分发的作用。

\item\textbf{网络或应用层基于感知结果进行调度、规划和控制}

在获得环境状态信息后，网络侧或应用侧可以进一步开展资源调度、路径规划、行为决策、风险判断和控制优化等工作。例如，边缘计算节点可依据感知结果进行实时分析，网络控制面可根据目标状态调整资源分配策略，应用平台也可结合业务规则生成具体控制动作。这样，感知信息便不再停留在“被采集”的层面，而是直接转化为可执行的业务决策依据。

\item\textbf{执行动作后再通过感知进行反馈修正}

在控制命令被执行之后，环境和目标状态往往会发生新的变化。ISAC可以通过持续感知对执行效果进行验证，并将结果再次反馈给网络或应用系统，用于进一步调整控制策略。这就形成了“感知—传输—决策—执行—再感知”的动态循环过程，使业务系统能够不断修正偏差、优化动作并提升整体运行效率。
\end{enumerate}

因此，ISAC形成的闭环能力并不是简单地把通信与感知功能叠加在一起，而是通过二者深度融合，使网络能够从环境认知出发，贯通状态获取、信息共享、智能决策与反馈控制全过程，为复杂业务提供连续、实时和自适应的支撑能力.

\subsubsection{The significance of business loop capability for typical applications}
\label{subsubsec:ch1-isac-core value-closed loop-applications}

ISAC的业务闭环能力对于未来典型应用场景具有重要意义，其价值不仅体现在功能增强上，更体现在支撑业务实时运行、提升系统自治水平和降低协同成本等方面,具体如下：

\begin{itemize}
\item 在车路协同场景中，环境感知是驾驶决策的重要基础。传统车联网主要解决车辆与车辆、车辆与路侧基础设施之间的信息交换问题，但车辆能否安全通行、是否需要减速避让、是否存在盲区目标等，往往仍依赖车载传感器或独立路侧感知设施。ISAC引入业务闭环后，网络能够更直接地参与道路环境感知与状态共享，帮助车辆及时获取前方目标、道路障碍、路口态势和交通流变化等信息，并据此支撑驾驶辅助、路径调整和协同控制，从而提升行车安全性与交通运行效率。

\item 在低空经济场景中，目标探测、轨迹管理和空域协同对闭环能力要求尤为突出。无人机配送、巡检、测绘和应急飞行等任务需要网络持续掌握飞行器位置、速度、航迹和周边空域动态，同时还要支持飞行指令下发、冲突规避和任务调整。ISAC所形成的闭环能力可以将目标探测、状态传输、空域管理决策和飞行反馈联结起来，使网络在低空目标识别、轨迹监管和协同调度中发挥更主动的作用，为低空业务的大规模、规范化运行提供基础支撑。

\item 在工业自动化场景中，状态监测、异常识别和控制反馈构成了生产过程中的核心闭环。现代工业系统要求网络不仅连接设备，还要理解设备运行状态、工艺变化和环境扰动，并在发现异常后迅速触发控制反馈。ISAC可以帮助网络实现对设备位置、运动状态、振动变化或作业区域动态的实时感知，并将这些信息用于生产调度、异常告警和控制优化。这样一来，网络就不再只是工业数据的传输通道，而成为支撑感知、控制和反馈联动的重要组成部分。

\item 在公共安全场景中，监测、预警和联动响应高度依赖业务闭环能力。无论是重点区域安防、人员聚集监测，还是灾害预警和突发事件处置，都需要系统及时发现异常、快速共享信息并推动相关部门协同行动。ISAC通过融合感知与通信能力，可以更高效地支撑目标监测、态势判断、风险预警和多部门联动响应。在此过程中，网络不仅承担事件信息上报和指令下发功能，还能够直接参与态势感知和反馈更新，从而提升公共安全系统的实时性和整体响应能力。   
\end{itemize}

总体来看，在这些典型应用中，业务闭环能力的意义在于使网络能够更深入地嵌入业务流程本身，帮助系统实现从“看见环境”到“理解环境”、从“传递信息”到“驱动决策”、从“执行动作”到“持续优化”的能力升级。

\subsubsection{The business closed-loop capability is the key for ISAC to achieve its 6G inherent capabilities}
\label{subsubsec:ch1-isac-core value-closed loop-6g}
从技术演进和网络体系升级的角度看，业务闭环能力是ISAC走向6G内生能力的关键。对于ISAC而言，只有当感知能力能够真正参与业务闭环，并在网络运行和应用支撑中发挥持续作用时，通信与感知的融合才具有体系化意义。

首先，业务闭环能力使感知不再只是附属测量功能，而成为网络服务能力的重要组成部分。在传统观念中，感知往往被视为一种补充能力，用于提供测距、定位、探测或环境观测等附加信息。但在ISAC框架下，感知结果直接影响网络调度、资源分配、业务控制和应用执行，其作用已经超出单纯“测量”的范畴，转而成为支撑网络智能化运行和业务动态优化的基础能力。也就是说，感知不只是“网络可选的附加功能”，而是网络提供高级服务时必须具备的关键组成部分。

其次，业务闭环能力为后续“ISAC+”以及“通感算控一体化”埋下了逻辑基础。未来6G网络的发展方向之一，是推动通信、感知、计算和控制等能力进一步协同，实现跨层联动和系统级优化。在这一过程中，ISAC如果仅停留在信号层或物理层的一体化设计，其作用仍然是局部的；而一旦ISAC具备完整的业务闭环能力，就能够自然向更高层次的融合演进。例如，感知结果可以触发边缘计算任务，计算结果可以指导控制策略，控制执行后的反馈又可以通过感知机制进行验证，从而形成更加完整的“通感算控一体化”体系。

因此，业务闭环能力不仅体现了ISAC相对于传统通信系统的功能优势，也决定了其能否真正成为6G网络中的基础能力和核心服务能力。只有当ISAC能够支撑业务全流程运行，嵌入感知、传输、决策与反馈的连续循环中，它才真正具备从技术概念走向网络内生能力的条件。

\subsection{Reflect questions}
\label{subsec:ch1.2-reflect-questions}
一、概念理解题
1. 为什么说ISAC不是“可有可无的功能叠加”，而是一种系统能力的重构？请结合“独立通信系统”“独立感知系统”和“融合系统”的差异进行说明。

2. 本节指出，ISAC的核心价值不是单点性能提升，而是系统级收益重构。请解释“系统级收益重构”的含义，并说明它与“某一项指标更优”之间的本质区别。

3. 为什么说传统通信系统本质上是“对物理世界无感”的系统？这种局限在5G以前可能并不突出，但在6G时代会变得越来越明显，请结合应用场景说明原因。

二、分析比较题
4. 请从频谱占用方式、硬件建设、部署流程和业务支撑能力四个方面，对比以下两种方案的差异：

（1）通信系统与感知系统独立部署；

（2）采用ISAC的一体化部署。

并分析哪一种方案更适合面向6G的大规模新型场景。

5. 本节从频谱效率、硬件成本、部署效率和业务闭环能力四个方面分析了ISAC的核心价值。请说明这四种价值之间是否彼此独立。如果不是，请分析它们之间可能存在的内在联系。

6. 频谱效率提升并不意味着“任何场景下ISAC都一定优于独立系统”。请结合以下因素，分析ISAC频谱增益的适用条件与边界：

双功能冲突
互扰与自干扰
目标分辨率要求
通信链路质量要求
三、深入思考题
7. 假设某城市计划建设车路协同基础设施，有两种方案：

方案A：分别建设通信基站和路侧雷达；
方案B：在通信基站上升级ISAC能力。
请从CAPEX、OPEX、部署周期、维护复杂度和后续业务扩展能力五个角度，对两种方案进行比较，并给出你的建议。
8. 本节提到ISAC通过共享天线阵列、射频链路、基带处理平台等硬件，能够降低系统总成本。请思考：

如果通信与感知对动态范围、线性度、时钟精度、收发隔离等要求不同，那么硬件共享为何仍然可能是划算的？

请尝试从“系统总成本”而不是“单模块最优性能”的角度进行论述。

9. 请围绕“部署效率”这一价值，回答以下问题：

为什么说ISAC特别适合低空经济、工业互联网、公共安全等新兴场景？

如果这些场景完全依赖独立感知基础设施，部署上会遇到哪些现实障碍？

10. 本节提出，ISAC的最大价值之一在于形成“感知—传输—决策—反馈”的业务闭环。请选择以下任一场景进行分析：

车路协同
低空无人机管理
智能工厂
公共安全监测
要求说明：

（1）感知环节获取什么信息；

（2）通信环节传输什么内容；

（3）谁来做决策；

（4）反馈如何作用于下一轮系统行为。

四、综合论述题
11. 有一种观点认为：“只要通信系统外接一些雷达或传感器，也能实现环境感知，因此没有必要发展ISAC。”

请你结合本节内容，对这一观点进行评述。要求至少从以下三个方面展开：

系统成本
部署效率
业务闭环能力
12. 请论述为什么说业务闭环能力是ISAC走向6G内生能力的关键，而不仅仅是一个附加优势。

提示：可以从网络角色变化、智能化需求、数字孪生、通感算控一体化等角度展开。

13. 假设未来6G网络已经普遍具备ISAC能力，那么运营商的角色会发生什么变化？

它是否仍然只是“连接服务提供者”，还是会进一步变成“环境信息服务提供者”甚至“决策支撑平台”？请谈谈你的理解。

五、开放探索题
14. 如果你是一名系统设计者，需要为某一类场景评估ISAC的实际价值，你会从哪些维度建立评估指标体系？

请至少给出5个评价维度，并说明每个维度对应反映的是ISAC的哪一类核心价值。

15. 请思考：在未来一段时间内，ISAC最可能首先在哪些行业实现规模化落地？

是车联网、低空经济、工业自动化、智慧城市，还是其他场景？

请结合“频谱效率、硬件成本、部署效率、业务闭环能力”四个维度给出理由。

16. 本节强调ISAC的价值是“系统级收益重构”。那么在实际商用中，是否可能出现这样的情况：

某一项单独指标（例如通信速率或感知精度）略有下降，但整个系统反而更有价值？

请举例说明，并讨论这是否会影响ISAC的产业接受度。